\chapter{The Gopher Protocol}
\label{chap:gopher-protocol}

\section{Historical Context}

Gopher, defined in RFC~1436 (1991), predates the World Wide Web's mainstream adoption and was, for a period in the early 1990s, a leading protocol for distributing hierarchically-organized documents over the Internet. Its design goals were simplicity and low implementation cost, both of which make it an excellent teaching protocol today.

\section{Protocol Overview}

A Gopher exchange is a single request followed by a single response, after which the connection is closed:

\begin{enumerate}
    \item The client opens a TCP connection to the server (traditionally port 70; this project uses port 7070 to avoid requiring root privileges, since ports below 1024 are typically restricted to privileged processes on Linux).
    \item The client sends exactly one line of text, called a \textbf{selector}, terminated by \texttt{\textbackslash r\textbackslash n}. An empty selector (i.e.\ just \texttt{\textbackslash r\textbackslash n}) requests the server's root menu.
    \item The server sends back the corresponding content as a raw byte stream.
    \item The server closes the connection. This closure \emph{is} the end-of-message signal --- there is no length field or other explicit terminator required by the protocol (though menus conventionally end with a line containing only a period, as described below).
\end{enumerate}

\section{The Menu Format}

When a selector refers to a directory-like resource rather than a plain document, the response is a \textbf{menu}: a list of further selectors the client can follow. Each line of a menu has the format:

\begin{center}
\texttt{<type><display\_text>\textbackslash t<selector>\textbackslash t<host>\textbackslash t<port>\textbackslash r\textbackslash n}
\end{center}

\begin{itemize}
    \item \textbf{type} is a single character indicating what kind of resource this entry is. This implementation uses \texttt{0} for a plain text file, \texttt{1} for a submenu (directory), and \texttt{3} for an error.
    \item \textbf{display\_text} is the human-readable label a client should show to the user.
    \item \textbf{selector} is the string the client should send back if it chooses this entry.
    \item \textbf{host} and \textbf{port} tell the client where to connect to fetch this entry (allowing a single menu to reference resources on entirely different servers).
\end{itemize}

Fields within a line are separated by the \textbf{tab} character (\texttt{\textbackslash t}), not spaces --- a detail that is easy to overlook but essential for compliant parsing. A menu response is conventionally terminated by a line containing a single period (\texttt{.\textbackslash r\textbackslash n}).

\section{Comparison Table: Gopher's Minimalism}

\begin{center}
\begin{tabularx}{\textwidth}{l X X}
\toprule
\textbf{Aspect} & \textbf{Gopher} & \textbf{HTTP (for comparison)} \\
\midrule
Request format & One line (the selector) & Request line + headers + blank line \\
Response format & Raw content & Status line + headers + blank line + body \\
End-of-message signal & Connection closure & \inlinecode{Content-Length} header \\
Status/error reporting & None (type \texttt{3} lines by convention) & Numeric status codes (200, 404, \ldots) \\
Methods & None --- always ``fetch this selector'' & GET, POST, PUT, DELETE, \ldots \\
\bottomrule
\end{tabularx}
\end{center}

This comparison is expanded further once HTTP itself has been covered in Chapter~\ref{chap:http-protocol}.
